\section{Operational Forensic Applications and Targets (RQ2)}
\label{sec:use_cases}

To answer RQ2, this section synthesizes empirical implementations across all six subdomains of the core corpus ($N=44$). Rather than operating as universal autonomous investigators, language models cluster around five DFIR operational bottlenecks (evidence triage and reporting, incident response and log correlation, malware analysis, volatile memory forensics, and covert artifact inspection) alongside a rapidly expanding inverse discipline: digital forensics \textit{of} LLM and agentic systems themselves.

\subsection{Evidence Triage, Reporting, and Benchmark Datasets}
Investigative triage, reporting, and dataset construction represent the largest operational cluster, accounting for 40.9\% ($n=18$) of the validated corpus. Examiners devote substantial effort to synthesizing fragmented artifacts into structured judicial documentation. Michelet and Breitinger~\cite{michelet2024_chatgpt} conducted a foundational experiment examining canonical forensic report structures and benchmarking ChatGPT (GPT-3.5) against a local open-weights baseline (Llama-2-13B) for assisted report drafting. Their evaluation showed that models compress multi-source case details into readable drafts, yet both commercial and local architectures exhibited non-negligible hallucinations and factual inaccuracies unsuitable for unverified court submission.

Subsequent engineering prioritized domain-specific adaptation and on-premise execution. Sharma et al.~\cite{sharma2025_forensicllm} introduced ForensicLLM, a 4-bit quantized LLaMA-3.1-8B architecture fine-tuned on digital forensics Q\&A samples so seized evidentiary material never traverses external networks. Michelet et al.~\cite{michelet2025_finetuning} formalized practical recommendations for instruction-tuning LLMs on forensic tasks, demonstrating through a chat summarization case study that domain fine-tuning improves structured extraction reliability over generic zero-shot prompting.

Complementary architectures couple language models with formal knowledge graphs or human-in-the-loop verification to anchor semantic outputs. Zhou et al.~\cite{zhou2025_forensic} proposed Forensic Intelligence Graphs to extract structured entity relationships from heterogeneous digital evidence, while Khatiwala et al.~\cite{khatiwala2025_evaluating} integrated LLM-refined artifact extraction with a Digital Forensic Knowledge Graph (DFKG), validating artifact traceability on a 13~GB Capture-the-Flag forensic disk image comprising 61 applications, 2,864 databases, and 5,870 tables. At the evidence-intake layer, Yoo and Jeong~\cite{yoo2026_scpllmscp} applied LLM keyword augmentation to improve title-driven file selection, Johnson et al.~\cite{johnson2026_aiassisted} evaluated multilingual online gaming chat triage, and Ko et al.~\cite{ko2026_humansupervised} validated human-supervised multi-model identifier extraction (Claude Sonnet 4.5, GPT-5.2, Gemini 3 Pro) across 273 California DFPI cryptocurrency-fraud ("pig butchering") complaints against dual-rater ground truth ($\kappa=0.969\text{--}0.988$), achieving $F_1=1.000$ on wallet addresses, $F_1=0.911\text{--}0.918$ on URLs, and $F_1=0.865\text{--}0.904$ on loss amounts (vs.\ $F_1 \le 0.381$ for regex/spaCy) while routing 31.1\% ($82/264$) of actionable cases to prioritized human review. In a domain-adjacent pilot study spanning photographic crime-scene triage, Farber~\cite{farber2025_scpaiscp} evaluated GPT-4, Claude, and Gemini across 30 crime-scene images scored by 10 forensic experts, concluding that frontier models serve strictly as preliminary screening aids rather than autonomous analysts. To establish reproducible triage evaluation, Chakir and Lahsen-Cherif~\cite{chakir2025_forensicsdata} introduced ForensicsData, a public benchmark containing over 5,000 Question-Context-Answer triplets derived from real malware analysis reports, while empirical law-enforcement case studies~\cite{kim2025_digital} and early exploratory frameworks~\cite{xu2024_transforming, xu2025_enhancing, wickramasekara2025_exploring} document practical integration barriers.

\subsection{Incident Response, Logs, and Agentic SOC Orchestration}
Incident response, log interpretation, and Security Operations Center (SOC) orchestration constitute 20.5\% ($n=9$) of the corpus. Enterprise intrusions generate high-velocity telemetry across Windows Event Logs, cloud audit trails, and network sensors that overwhelm manual correlation.

Recent studies shift from monolithic single-prompt scripts toward modular agentic orchestration and specialized lightweight models. Liu~\cite{liu2025_autobnb} proposed AutoBnB, a multi-agent incident response architecture partitioning log triage, analysis, and containment reasoning across specialized LLM roles. Similarly, {\"O}zg{\"o}zler et al.~\cite{zgzler2026_eventdriven} designed an Event-Driven Agentic SOC (ED-ASOC) combining supervisor-coordinated LLM agents with SOAR workflows, while Matsuda et al.~\cite{matsuda2025_enhancing} integrated Sigma detection rules, Model Context Protocol interfaces, and operational policies into LLM-assisted response. In cloud environments, Alharthi and Yasaei~\cite{alharthi2025_llmpowered} automated cloud forensic log investigation, and George et al.~\cite{george2025_llm} fused Cyber Threat Intelligence (CTI) knowledge graphs with LLMs to map observed telemetry onto adversary tactics. Extending agentic root-cause analysis to adjacent distributed software operations, Tian et al.~\cite{tian2026_gala} introduced GALA, combining graph-augmented agents for microservice incident triage.

Addressing hallucination and inference cost directly, Hammar et al.~\cite{hammar2026_incident, hammar2026_a} developed fine-tuned lightweight LLM decision-support systems for incident response planning evaluated across testbed and synthetic attack scenarios, demonstrating reduced hallucination rates compared to unconstrained frontier prompting. For forensic timeline evaluation specifically, Studiawan et al.~\cite{studiawan2025_towards} designed a standardized quantitative methodology and dataset modeled on the NIST CFTT framework~\cite{nist_cftt}.

A critical failure mode in automated log and timeline synthesis is narrative coherence bias. Autoregressive models optimize token fluency, compelling disparate log records into a smooth causal storyline. When raw timelines exhibit clock skew, NTP drift, or deliberate anti-forensic timestomping (such as discrepancies between NTFS \texttt{\$STANDARD\_INFORMATION} and \texttt{\$FILE\_NAME} timestamps), unconstrained language models silently smooth away anomalies. Chronological smoothing destroys dispositive evidence. Sound forensic engineering therefore requires deterministic engines such as \texttt{plaso} or \texttt{log2timeline} to fix immutable UTC event sequences before any LLM performs semantic annotation.

\subsection{Malware Analysis and Reverse Engineering}
Malware triage and reverse engineering represent 13.6\% ($n=6$) of the corpus. Modern malicious binaries evade static signature matching via packing, polymorphism, and obfuscated control flow.

Empirical evaluations demonstrate complementary language-model paradigms across behavioral and static layers. Clark and Chen~\cite{clark2025_large} fine-tuned a pretrained RoBERTa transformer on structured Windows application-layer API and system-call traces, capturing behavioral semantics between benign and malicious execution; their 512-token \texttt{prob\_mean} aggregation configuration achieved 92\% recall and an $F_1$ score of 0.94 at the application level. Complementing behavioral analysis, Meskova et al.~\cite{meskova2026_aidriven} applied transformer detection to real-time malware threat analysis, while Jelodar et al.~\cite{jelodar2026_large} synthesized broader LLM capabilities across source-code inspection, vulnerability analysis, and reverse engineering.

To overcome context-window limits on large binaries, recent architectures fuse LLMs with structural program graphs or multi-agent decomposition. Sava et al.~\cite{sava2026_triagecpg} introduced TRIAGE-CPG, pairing Code Property Graphs (CPGs) with eight closed-source and open-weight models evaluated across 1,000 Windows PE binaries (600-sample test split); inspecting roughly 40 of a mean 1,000 functions per binary, the top model reached 0.805 strict balanced accuracy and 0.855 triage balanced accuracy (flagging 65\% of malware, escalating 12\%, and clearing 94\% of benign binaries). Zheng et al.~\cite{zheng2025_avagent} developed AV-Agent, a bottom-up interpretable malware classifier producing auditable attribution chains, and Ozyurt et al.~\cite{ozyurt2026_llmorchestrated} orchestrated parallel three-agent Android APK inspection combining Gemini 2.5 Flash, DeepSeek V3, and GPT-4.1-mini via majority vote.

\subsection{Volatile Memory Forensics (LLM-Assisted)}
LLM-assisted volatile memory forensics accounts for 6.8\% ($n=3$) of the corpus (distinct from memory acquisition \textit{of} LLM applications discussed below). Interpreting raw RAM dumps via Volatility~\cite{case2014volatility} requires deep kernel expertise to separate injected code, unlinked libraries, and hidden network sockets from benign OS state.

Oh et al.~\cite{oh2024_volgpt} introduced volGPT, automating ransomware process triage from structured Volatility outputs to detect process masquerading and volatile artifacts at substantially lower triage latency than manual inspection. Addressing analyzer performance and memory safety, Gharaibeh and Baggili~\cite{gharaibeh2026_rusting} conducted an empirical LLM-assisted migration of 38 Volatility3 Linux plugins from Python to Rust (29,028 lines of code over 68 hours, consuming 900 million tokens at \$2,962); the Rust port achieved a 2.50$\times$ median speedup (up to 210$\times$ on structure traversal), completed 417/417 executions (vs.\ 11.3\% Python runtime failures), and matched canonical output hashes across 68.9\% of comparable runs. Targeting financial privacy in RAM examinations, Kao et al.~\cite{kao2026_a} proposed a local hierarchical LLaMA-3.1-8B Tri-Layer RAG-Decider framework for privacy-preserving cryptocurrency-wallet memory forensics, achieving a 0.84 BERT $F_1$ score across 100 wallet-forensics queries and performing within 0.3\% of commercial ChatGPT-4o overall while surpassing it on reasoning.

\subsection{Digital Forensics of LLM and Agentic Systems}
A rapidly emerging inverse discipline, accounting for 11.4\% ($n=5$) of the corpus, treats LLM runtimes and agentic frameworks not as analytical tools, but as investigative targets generating novel digital evidence. As local inference engines and autonomous agents execute code, browse files, and invoke tools on suspect endpoints, examiners must recover what prompts were issued and which actions executed.

Chernyshev et al.~\cite{chernyshev2023_towards} initiated this axis by analyzing LLM invocation logs to establish forensic traceability of model usage. Walker et al.~\cite{walker2024_forensic} systematically cataloged persistent disk and execution artifacts generated by Microsoft's multi-agent AutoGen platform. At the volatile memory layer, Xu et al.~\cite{xu2026_recovering} demonstrated at IEEE S\&P that conversational histories and cached inference contexts from local mobile LLMs can be recovered and rehosted directly from smartphone RAM dumps. Similarly, Satter et al.~\cite{satter2026_with} demonstrated memory forensic reconstruction of Model Context Protocol (MCP) activity in agentic LLM systems even when application logging is disabled or absent. Synthesizing this emerging domain, Scanlon~\cite{scanlon2026_investigation} formalized digital forensic readiness, evidentiary artifact challenges, and regulatory gaps when investigating proprietary and local AI systems.

\subsection{Admissibility, Explainability, and Covert Artifacts}
The remaining 6.8\% ($n=3$) of the corpus directly targets judicial admissibility, neuro-symbolic explainability, and covert binary artifacts. Trivedi~\cite{trivedi2026_when} empirically tested examiner bias, prompt framing, and evidentiary admissibility across 2,700 trials (analyzed in depth in Section~\ref{sec:custody_crypto}). Tiwari and Nene~\cite{tiwari2026_xforensicnet} introduced X-ForensicNet, a fully offline multi-agent Windows forensics framework coupling local LLM inference, RAG, and neuro-symbolic validation, reporting a 2.8$\times$ investigative speedup and 96.3\% accuracy with explainable reasoning trails. Beyond textual artifacts, Sudheer et al.~\cite{sudheer2026_least} explored LLM-assisted detection and contextual analysis of Least Significant Bit (LSB) image steganography, showing that while generative models cannot replace statistical steganalysis detectors, they assist analysts in triaging extracted covert payloads.

Table~\ref{tab:taxonomy_summary} synthesizes all six operational subdomains, highlighting verified exemplar studies, dominant architectures, and primary forensic safeguards.

\begin{table*}[!t]
\caption{Taxonomic Synthesis of Core Literature Across All Six Mutually Exclusive Subdomains ($N=44$)}
\label{tab:taxonomy_summary}
\centering
\scriptsize
\setlength{\tabcolsep}{3pt}
\begin{tabular}{p{2.2cm}cp{3.7cm}p{2.1cm}cp{3.9cm}}
\toprule
\textbf{Subdomain} & \textbf{Corpus} & \textbf{Verified Exemplar Studies} & \textbf{Dominant Paradigm} & \textbf{RQs} & \textbf{Forensic Relevance \& Safeguards} \\
\midrule
Evidence Triage, Reporting \& Datasets & 18 (40.9\%) & ForensicLLM~\cite{sharma2025_forensicllm}, Reporting~\cite{michelet2024_chatgpt}, DFKG~\cite{khatiwala2025_evaluating}, ForensicsData~\cite{chakir2025_forensicsdata} & Cloud APIs / Local Fine-Tuned & RQ2, RQ3, RQ4 & Case-note synthesis, graph grounding, and benchmark corpora; requires air-gapped local weights and prompt-injection defense. \\
\addlinespace
Incident Response, Logs \& Agentic SOC & 9 (20.5\%) & AutoBnB~\cite{liu2025_autobnb}, Lightweight IR~\cite{hammar2026_incident}, Cloud Logs~\cite{alharthi2025_llmpowered}, Timelines~\cite{studiawan2025_towards} & Multi-Agent / Lightweight LLMs & RQ2, RQ4 & Distributed telemetry correlation and response planning; requires immutable \texttt{plaso} timelines to block narrative smoothing. \\
\addlinespace
Malware Analysis \& Reverse Eng. & 6 (13.6\%) & RoBERTa API Calls~\cite{clark2025_large}, TRIAGE-CPG~\cite{sava2026_triagecpg}, AV-Agent~\cite{zheng2025_avagent}, APK Agents~\cite{ozyurt2026_llmorchestrated} & Fine-Tuned / Graph + Agentic & RQ2, RQ4 & System-call classification and CPG triage; vulnerable to adversarial call/code perturbation and obfuscation. \\
\addlinespace
Forensics of LLM \& Agentic Systems & 5 (11.4\%) & AutoGen~\cite{walker2024_forensic}, Mobile RAM~\cite{xu2026_recovering}, MCP Memory~\cite{satter2026_with}, AI Readiness~\cite{scanlon2026_investigation} & Target Subject (Disk / RAM / Logs) & RQ1, RQ2, RQ3 & Recovers prompts, tool invocations, and KV/context buffers from suspect LLM runtimes under ISO/IEC 27037/27042. \\
\addlinespace
Volatile Memory Forensics (Assisted) & 3 (6.8\%) & volGPT~\cite{oh2024_volgpt}, Rusting Volatility~\cite{gharaibeh2026_rusting}, Wallet RAM~\cite{kao2026_a} & Local / Cloud + Volatility & RQ2, RQ3 & Ransomware triage, Rust plugin migration, and privacy-preserving wallet memory analysis; mandates isolated RAM scratch. \\
\addlinespace
Admissibility, XAI \& Covert Artifacts & 3 (6.8\%) & Premise Bias~\cite{trivedi2026_when}, X-ForensicNet~\cite{tiwari2026_xforensicnet}, LSB Stego~\cite{sudheer2026_least} & Offline Neuro-Symbolic / Audit & RQ3, RQ4 & FRE 702(d) premise auditing, deterministic neuro-symbolic grounding, and covert payload triage. \\
\midrule
\textbf{Total Corpus} & \textbf{44 (100.0\%)} & \multicolumn{4}{p{10.8cm}}{\textit{Mutually exclusive partition across Axis A ($n=39$, LLMs for DFIR) and Axis B ($n=5$, Forensics of LLM Systems).}} \\
\bottomrule
\end{tabular}
\end{table*}

